\section{Quasinormal poles and the two-reflection sector in \texorpdfstring{AdS\(_7\)}{AdS7}}
\label{sec:einstein-poles}

At fixed scalar mass, the terms dropped in \eqref{eq:bessel-core} are
of the same order as its first mass correction. To combine them,
put \(x=\int_r^\infty\dd r'/f(r')\), \(D=d-1\),
\(\psi=r^{D/2}\phi\) and \(\nu=\sqrt{d^2/4+M^2}\) in
\eqref{eq:wave}. The exact potential is
\begin{equation}
 -\partial_x^2\psi+V\psi=\omega^2\psi,\qquad
 V=f\left[\nu^2-\frac14+\frac{D^2}{4r^d}\right].
 \label{eq:einstein-potential}
\end{equation}
Choose the horizon detour for which the singularity is at
\(x=\mathcal Z=T_*^*/2\). With \(s=x-\mathcal Z\) near zero,
\begin{equation}
 s=\frac{r^D}{D}\left(1+a r^d+b r^{2d}+O(r^{3d})\right),
 \qquad a=\frac{D}{2d-1},\quad b=\frac{D}{3d-1}.
\end{equation}
Subtracting the inverse-square term gives
\begin{equation}
 \begin{aligned}
 V+\frac{1}{4s^2}
 &=-A_\nu r^{-(d-2)}+B_\nu^{\rm E}r^2+O(r^{d+2}),\\
 A_\nu&=\nu^2-\frac{d^2}{4(2d-1)},\\
 B_\nu^{\rm E}&=\nu^2-\frac14+\frac{D^2}{4}(3a^2-2b).
 \end{aligned}
 \label{eq:einstein-shifts}
\end{equation}
At \(\nu_*={d}/[{2\sqrt{2d-1}}]\), the first Puiseux term
vanishes. The residual potential begins at
\(B_{\nu_*}^{\rm E}(D s)^{2p-2}\), with \(p=d/D\) as in
section~\ref{sec:connection}.

Use \(z=\ii\omega s\) and \(\psi=\sqrt z F(z)\). The linear
correction to the inner equation is then a term
\(g_P z^{P-1}F\) in \eqref{eq:bessel-core}. For generic \(\nu\)
take \(P=p\) and \(g_p=A_\nu D^{p-2}(\ii\omega)^{-p}\).
At \(\nu=\nu_*\) take \(P=2p\) and
\(g_{2p}=-B_{\nu_*}^{\rm E}D^{2p-2}(\ii\omega)^{-2p}\).
The coefficient \(\sigma(P)\) of \(g_P\) in the three-ray
multiplier is
\begin{equation}
 \sigma(P)=2[\cos(\pi P)-1]
 \left[-\frac{\tan(\pi P/2)}{\pi}\,
       \frac{2^{P-3}\Gamma(P/2)^4}{\Gamma(P)}\right].
 \label{eq:einstein-stokes-slope}
\end{equation}
For \(P=2p\), this use of \eqref{eq:bessel-moments} is at the
level of the linear perturbation of \(K_0\). The integral with
\(K_0^2\) converges and the integral with \(I_0K_0\) uses the
same finite-part subtraction \(R^{P-1}/[2(P-1)]\).
After subtracting \(1/(2t)\) from \(I_0(t)K_0(t)\), the large
\(t\) integrand is \(O(t^{P-4})\), which is integrable for
\(P<3\). Rotating the linear term gives \(e^{\pm\ii\pi P}\)
on the adjacent rays. This establishes \eqref{eq:einstein-stokes-slope}
without imposing a finite-coupling boundary condition at
\(P=2p>2\).

For \(d=6\), use the exterior transport in eqs.~(14.40) and (15.26)
of ref.~\cite{Jia:2026} to connect these core data to boundary observables. In its
positive lateral convention, the primitive transport factor is
\begin{equation}
 q(\omega)R(\omega),\qquad
 q(\omega)=e^{2\ii\mathcal Z\omega-\ii\pi\nu},\qquad
 R(\omega)=-2-\sigma(P)g_P+O(\omega^{-P-p}) .
 \label{eq:einstein-primitive-transport}
\end{equation}
The boundary phase in \(q\) comes from the normalizable Bessel
solution. The nonspiralling Stokes crossing in AdS\(_7\) contributes
at the thermal scale \(e^{-\beta\omega}\). In the pole
condition this crossing gives an exponentially small correction
on the positive-real-part quasinormal ray. The algebraic pole
positions follow from \(1-qR=0\). With the overtone index fixed
by the leading expression, they are
\begin{align}
 \omega_n^{(0)}&=\frac{\pi[n+(\nu-1)/2]+\tfrac{\ii}{2}\log2}
                         {\mathcal Z},\notag\\
 \mathcal Z(\omega_n-\omega_n^{(0)})
   &=\frac{\ii}{4}\sigma(P)g_P(\omega_n^{(0)})
       +O\big((\omega_n^{(0)})^{-P-p}\big).
 \label{eq:einstein-qnm-shift}
\end{align}
The phase branch of \((\ii\omega)^{-P}\) is continued from
\(\omega>0\) towards \(\arg\omega=-\pi/6\). At
\(\nu_*=3/\sqrt{11}\), the coefficient of
\((\omega_n^{(0)})^{-12/5}\) in
\(\mathcal Z(\omega_n-\omega_n^{(0)})\) is
\(0.0258400+0.0355657\ii\).

\begin{table}[t]
 \centering\small
 \begin{tabular}{ccccc}
 \hline
 \(\nu\) & \(n\) & \(\omega_n\) & Scaled shift & Limit \\
 \hline
 0.500000 & 5 & $12.514512-6.840325\ii$ & $-0.567148-0.002496\ii$ & -0.568182 \\
0.500000 & 15 & $38.489210-21.838852\ii$ & $-0.567957-0.000633\ii$ & -0.568182 \\
0.500000 & 30 & $77.459154-44.338583\ii$ & $-0.568089-0.000273\ii$ & -0.568182 \\
0.904534 & 5 & $13.031844-7.141721\ii$ & $1.008068-0.057591\ii$ & 1.000000 \\
0.904534 & 15 & $39.012606-22.141782\ii$ & $1.003870-0.014267\ii$ & 1.000000 \\
0.904534 & 30 & $77.983749-44.641786\ii$ & $1.001835-0.006102\ii$ & 1.000000 \\
1.250000 & 5 & $13.470821-7.398442\ii$ & $0.746673-0.006911\ii$ & 0.744318 \\
1.250000 & 15 & $39.458697-22.400284\ii$ & $0.744944-0.001874\ii$ & 0.744318 \\
1.250000 & 30 & $78.431348-44.900630\ii$ & $0.744589-0.000819\ii$ & 0.744318 \\
\hline

 \end{tabular}
 \caption{Poles of the full scalar equation in planar Schwarzschild--AdS\(_7\).
 The scaled shift is the left-hand side of \eqref{eq:einstein-qnm-shift}
 divided by its leading prediction for \(\nu=\nu_*\). For the other
 masses, it is divided by that prediction per unit \(A_\nu\).}
 \label{tab:einstein-poles}
\end{table}

For table~\ref{tab:einstein-poles} we solve the full boundary-value
problem using \(u=1/r\). The series at \(u=0\) has leading power
\(u^{1/2+\nu}\) and the ingoing series at \(u=1\) has leading power
\((1-u)^{-\ii\omega/d}\). We match their logarithmic derivatives
at \(u=0.53\), then recompute the poles at \(u=0.61\) with more
series terms and higher precision. The two calculations agree to
within \(10^{-10}\) in \(\omega_n\). At the tuned mass, the complex
ratio of the measured shift to the \(\omega^{-12/5}\) prediction
is within \(0.7\%\) of one at \(n=30\).

\subsection{Time dependence of the primitive two-reflection sector}

For real positive \(\omega\), write \(2\mathcal Z=A+\ii\beta/2\)
with \(A=(\beta/2)\cot(\pi/6)\). The transfer formula of
ref.~\cite{Jia:2026} gives the primitive second harmonic
through the required order in \(\rho/\rho_{\rm vac}\) as
\(q^2 R^2\) and its complex conjugate. Indeed, expanding the
retarded ratio of eq.~(14.40) of ref.~\cite{Jia:2026} gives
\(1+(1-e^{2\ii\pi\nu})(qR+q^2R^2)\) through quadratic
order in the primitive symbol. Taking the imaginary part with the
vacuum phase \(e^{-\ii\pi\nu}\) gives \(2\operatorname{Re}(q^2R^2)\)
at this harmonic. At this exponential order the boundary
normalization in the median prescription and the mixed product
of the two reflections contribute at zero real time displacement.
The thermal factor that converts \(\rho\) to \(G^>\) also changes
the zero-displacement term at this order. The perturbative
boundary factor starts at relative order \(\omega^{-d}\).

At the tuned mass, write \(R=-2+r_2\omega^{-12/5}+\cdots\),
where \(r_2=\sigma(2p)B_{\nu_*}^{\rm E}D^{2p-2}e^{-\ii\pi p}\).
Using \(\rho_{\rm vac}=C_{\nu_*}\omega^{2\nu_*}\) for positive
frequency, and \(w=\ii(t-4\mathcal Z)\), the Laplace integral
of the primitive positive-displacement two-reflection sector
then has the singular terms
\begin{equation}
 G^>_{2,+}(t)\ \supset+ \frac{4 C_{\nu_*} e^{-2\ii\pi\nu_*}}{2\pi}
 \left[
 \frac{\Gamma(2\nu_*+1)}{w^{2\nu_*+1}}
 -r_2\frac{\Gamma(2\nu_*+1-12/5)}{w^{2\nu_*+1-12/5}}
 \right].
 \label{eq:einstein-two-time}
\end{equation}
The second term diverges with power \(2\nu_*+1-12/5\simeq0.4091\).
Its coefficient is fixed by the same local potential term that
controls the tuned pole shift. The time \(4\mathcal Z\) equals
\(\sqrt3\,\beta+\ii\beta\), whereas a geodesic with any fixed
number of local cycles approaches \(T_*\). Equation
\eqref{eq:einstein-two-time} is a statement about the
high-frequency sector. The exact sheet-resolved continuation of
the full Einstein correlator would also require controlling the
other frequency sectors.
